\section{A discovery journey that preserves refusal and uncertainty}
\label{exec13:discovery-journey}

A bank's request to discover why an offer changes twelve-month customer value
contains several operations. The information must precede the offer, the
candidate graph must retain its unresolved orientations, the requested effect
must be identified under a declared model, and the financial interpretation
must preserve the resulting uncertainty. A common command-line or library
interface is useful only if its callers can tell which of these operations
actually occurred.

The local reference service supports graph syntax review and a particular
supplied-DAG back-door calculation. Its \texttt{graph.discover} request remains
unsupported. A syntax response can preserve a CPDAG or PAG and its endpoint
marks without establishing that the graph is true or recognizing every valid
member of the represented class. The supplied-DAG identification method does
not thereby become a method for arbitrary graph classes, latent confounding,
temporal discovery or dependent sampling.

The extended local exercise froze 24 authored cases and 98 invocation slots
before executing them~\citep{exec13_discovery_journey}. Each case was called
through the library and CLI, initially and after restart or re-instantiation.
Two additional requests retained a completed request's key while changing its
target horizon. The elementary references were graph paths, factored
inequalities, exact timestamp order and explicit changes to request identity.
They support a deterministic behavior comparison, not a statistical discovery
accuracy or causal-coverage estimate.

\subsection{Identification and refusal in a concrete offer graph}

For the supplied graph $A\to Y$, an empty adjustment set satisfies the selected
back-door criterion. Add a pre-offer risk variable $R$ with
$A\leftarrow R\to Y$. If $R$ is observed and included in the adjustment
request, the method returns the corresponding conditional identification
expression. Omitting $R$ leaves the path open. Declaring it latent does not
make an observed adjustment possible. Conversely, requesting adjustment for
the mediator in $A\to R\to Y$ fails this particular criterion for the total
effect. Such failures must remain failures of the attempted method; they are
not proofs that every identification strategy is impossible.

The cases also require refusal when a declared premise is contradicted,
unsupported dispositions for CPDAG/PAG identification and dependent-month
sampling, and an invalid-request response for a directed cycle. Correctly
recognizing these boundaries is useful behavior. It does not supply the causal
effect that the bank originally wanted.

\subsection{Numerical sets, information clocks and immutable context}

Four deterministic sensitivity calculations check set serialization:
$1-x^2\le0$ gives two rays, $1+x^2\le0$ gives the empty set,
$-1-x^2\le0$ gives the whole line, and $(x-1)^2\le0$ gives a singleton.
The service retains those different topologies. The examples have no sampling
coverage probability; numerical edge cases outside these selected polynomials
still need an independent error analysis.

The observation audit admits the timestamp row satisfying
$\text{event}\le\text{available}<\text{decision}$. Availability exactly at
the decision, a missing timestamp, and an availability time preceding the event
are retained as failures. These declarations alone cannot establish that a
real source system recorded the correct timestamps.

Two artifact cases test a different failure mechanism. Changing the sign of
the target after importing a moment artifact must prevent its reuse under the
new scientific context. Corrupting the stored bytes must fail their identity
check. Keeping an idempotency key while changing the horizon must yield a
conflict. A client that receives such a failure has no new computed effect to
display. Its request correlation ID can change across a replay, while the
scientific content must remain equal.

\subsection{Observed behavior and the cost of this local exercise}

All 98 slots matched their frozen expectations, and all 24 case comparisons
retained the same scientific responses across the four route/restart calls.
There were 44 computed responses, 36 unsupported responses, four refusals,
four invalid requests and ten unavailable responses, including the two target
conflicts. The computed group includes capabilities, syntax, conditional
expressions, sensitivity sets and an audit; it is not a count of 44 estimated
causal effects. All groups remain in the denominator.

The run resumed after 40 and 80 recorded invocations, preserving both completed
prefixes exactly. The 98 invocations took approximately 41.99 seconds in this
local run, including worker/process launch, failed requests and replay. This
measurement excludes author effort, process startup before the runner's timer,
and pauses between runner segments. Separate active-runner time includes
artifact imports and journaling. Reported resident-memory readings are maxima
for individual parent/child processes, not concurrent process-tree memory.
These measurements neither establish bank throughput nor estimate bank ROI.

Restarting these calls is a narrower exercise than killing a multi-host
service mid-operation. A crash between a completed call and its journal write
can repeat that call with the same key; the exercise does not establish
exactly-once execution. It also runs no actual sibling consumer, MCP transport,
live language model or installed discovery engine. Those distinctions preserve
the meaning of both positive checks and remaining work.

\subsection{How the proposed consumer uses the result}

A consumer first reads capabilities and checks its clock, target, sampling law,
graph type and requested method. A syntax pass allows a graph to be displayed
with its marks intact. A conditional identification expression permits the
next estimation design to be considered. An unsupported or refused result
routes the question to another justified method or requests missing information.
It must not trigger an automatic point estimate from a language model.

A contract double can demonstrate these client decisions using typed local
responses. Acceptance of an actual workspace integration additionally needs
the owner's request and unit map, source clock, execution trace, failure and
restart results, and an explicit disposition. This is later evidence about
that integration. For the present proposal, the local exercise already makes
the required behavior concrete and reproducible while preserving the missing
discovery-to-effect capabilities.
\label{exec13:discovery-journey-end}
